\documentclass[oneside]{article}
\usepackage{marcoreport}

\title{Test Build}
\author{Marco Zhang}
\date{February 2026}

\begin{document}
\maketitle

\section{Where We At}

We have a consumption-savings problem with utility $u(c) = \log(c)$ and budget:
\[
a_{t+1} = (1+r) a_t - c_t
\]

\section{Why We Get Motivated}

The value function is:
\[
V(a_0) = \sum_{t=0}^{\infty} \beta^t \log(c_t^*)
\]
But this solves only for a specific $a_0$. We want a rule for any wealth level.

\section{What We Do}

Separate the first term and re-index:
\begin{align}
V(a_0) &= \log(c_0^*) + \beta \sum_{t=1}^{\infty} \beta^{t-1} \log(c_t^*) \\
       &= \log(c_0^*) + \beta V(a_1)
\end{align}

Define vectors $\vct{x} \in \RR^n$ and matrix $\mat{A}$. The expectation $\EE[c_t^*]$ and $\var(a_t)$ are finite. The optimal policy is:
\[
c_t^* = \argmax_{c} \left\{ \log(c) + \beta V((1+r)a_t - c) \right\}
\]

\begin{remarkbox}[Bellman Equation]
This recursive structure $V(a) = \max_c \{ u(c) + \beta V(a') \}$ is the Bellman equation. It reduces an infinite-horizon problem to a single-period optimization.
\end{remarkbox}

\section{Results and Next Steps}

\begin{theorem}
Under $\beta(1+r) < 1$, the value function $V$ is concave and the policy function $c^*(a)$ is increasing in $a$.
\end{theorem}

\begin{lemma}
The envelope condition gives $V'(a) = (1+r) u'(c^*(a))$.
\end{lemma}

\begin{definition}
A \textbf{stationary equilibrium} is a distribution $\mu$ over wealth such that $\mu$ is invariant under the optimal policy.
\end{definition}

\begin{itemize}
  \item Bellman equation established
  \item Next: solve via FOC + envelope theorem
  \item Then: characterize the stationary distribution
\end{itemize}

\begin{tcolorbox}[mynote, title=Note]
This is a test of the mynote tcolorbox style from marcoreport.sty.
\end{tcolorbox}

\end{document}
